% !TeX root = ../../Book2.tex
\section{Mihlin–Hörmander 定理}
\label{b2:sec:2202}

导数条件的作用是使每个频率环带上的核有足够快的衰减，再让环带总和满足奇异积分的抵消估计。证明须同时使用符号的有界性来建立 \(L^2\) 起点；核的大小条件本身不能替代这个起点。

\subsection{符号导数}
\label{b2:subsec:220201}

设 \(N\) 为整数。常用的尺度不变符号界为
\[
 |\partial^\alpha m(\xi)|\le A|\xi|^{-|\alpha|},
 \qquad 0\le|\alpha|\le N,\quad \xi\ne0.
\]
零阶包含 \(|m|\le A\)。例如 \(\xi_j/|\xi|\) 的各阶导数都是相应次数的齐次函数，在单位球面连续有界，因此满足任意有限阶条件。相反，\(|\xi|^{-\alpha}\) 在零频无界，不能成为 \(L^2\) 上的有界零阶乘子。

\ProseParagraphBreak
若 \(m(\xi)=b(\log|\xi|)\)，而 \(b\) 的有限阶导数有界，则链式法则给出同样的尺度界。符号不需要在零点有极限，甚至可以沿尺度持续振荡。这里对零点的忽略不影响 \(L^p\) 算子，却提醒我们不能把定理误述成“符号在闭空间上光滑”。

\subsection{二进分解}
\label{b2:subsec:220202}

使用\chapref{b2:ch:19}的环带分割 \(\varphi(2^{-j}\xi)\)，令
\[
 m_j(\xi)=m(\xi)\varphi(2^{-j}\xi),\qquad
 k_j(x)=2^{jd}h_j(2^jx),\quad
 h_j=\mathcal F^{-1}\bigl[m(2^j\,\cdot)\varphi\bigr].
\]
在固定环带上，归一化符号的有限阶导数一致有界。对 Fourier 积分分部积分，并分别选择绝对值最大的坐标，可得
\[
 |h_j(x)|+|\nabla h_j(x)|\le C A(1+|x|)^{-N}.
\]
梯度估计只是把符号乘上有界环带中的 \(2\pi\mathrm i\xi_\ell\)，仍使用相同阶数的符号导数；不需要额外的一阶 \(m\) 导数。该估计对所有 \(j\) 一致。

\ProseParagraphBreak
当 \(N>d+1\)，在 \(x\ne0\) 上绝对求和：
\[
 \sum_j|k_j(x)|\le CA|x|^{-d},\qquad
 \sum_j|\nabla k_j(x)|\le CA|x|^{-d-1}.
\]
例如第一式把 \(2^j|x|\le1\) 与 \(>1\) 分开，分别求等比级数
\(\sum2^{jd}\) 和 \(|x|^{-N}\sum2^{j(d-N)}\)。
第二式同理需要 \(N>d+1\)。这些级数在离原点的紧集上一致收敛，从而其和 \(K\) 是那里的一阶光滑函数。

\subsection{Mihlin 乘子定理}
\label{b2:subsec:220203}

下面采用\term[mihlin-chengzi-dingli]{Mihlin 乘子定理}[Mihlin multiplier theorem]的下述版本。

\begin{theorem}[Mihlin 乘子定理的有限导数版本]\label{b2:thm:22-mihlin}
设 \(m\in C^{d+2}(\mathbb R^d\setminus\{0\})\)，且对
\(|\alpha|\le d+2\) 满足上述尺度导数界。则
\[
 \|T_mf\|_p\le C_{d,p}A\|f\|_p,\qquad1<p<\infty.
\]
该导数数目是本证明使用的方便充分条件，并不声称最优。
\end{theorem}
\begin{proof}
Plancherel 给出 \(L^2\) 范数至多 \(A\)。有限环带和的符号在几乎处处频率上趋于 \(m\)，且一致有界，因此对应算子在 \(L^2\) 中趋于 \(T_m\)。
取紧支撑 \(f\in L^2\)，其支撑之外的任意紧集到支撑的距离为正。前述核级数在那里一致绝对可积，可交换求和和积分，得到
\(T_mf(x)=\int K(x-y)f(y)\,\mathrm dy\)。
更准确地，先在两个不交支撑的测试函数之间配对，再取 \(L^2\) 极限，即确认这就是给定算子的支撑外核，而不是另行定义的算子。

核的大小与梯度界满足\chapref{b2:ch:18}的标准 Calderón--Zygmund 条件，且卷积核的两变量差分由同一个梯度界控制。\cref{b2:thm:18-cz-lp}因而给出所有 \(1<p<\infty\) 上的界。对 \(L^2\cap L^p\) 的界先建立，再以该空间在 \(L^p\) 中的稠密性定义唯一延拓。
\end{proof}
证明实际上还给出弱型 \((1,1)\)，但通常没有强型 \((1,1)\) 或 \(L^\infty\to L^\infty\) 界。Hilbert 变换的区间输入产生对数尾和边界对数奇性，已分别说明这两个端点为何不能由定理自动补齐。

\ProseParagraphBreak
有限频率环带上的导数条件也可使用：若 \(m\) 支撑在固定环带并有 \(d+1\) 阶导数，则它的核属于 \(L^1\)，从而包括端点在内都可有界。真正造成端点困难的是无限尺度叠加，不能把单个环带的 Young 估计在没有可求和常数时直接相加。

\subsection{Hörmander 条件}
\label{b2:subsec:220204}

导数的逐点控制可改为环带上的 Sobolev 控制。固定 \(\eta\in C_c^\infty\)，在 \(\varphi\) 支撑邻域等于 \(1\)，令
\[
 B_s(m)=\sup_{j\in\mathbb Z}\|m(2^j\,\cdot)\eta\|_{H^s(\mathbb R^d)}.
\]
这里 \(H^s\) 的变量是频率变量本身；此处的 Sobolev 范数不是输入 \(f\) 的范数。
下面采用\term[hormander-chengzi-dingli]{Hörmander 乘子定理}[Hörmander multiplier theorem]的下述版本。

\begin{theorem}[Hörmander 环带条件]\label{b2:thm:22-hormander}
若 \(s>d/2\)、\(m\in L^\infty\) 且 \(B_s(m)<\infty\)，则 \(m\in M_p\)、\(1<p<\infty\)，算子范数不超过
\(C_{d,p,s}(\|m\|_\infty+B_s(m))\)。
\end{theorem}
\begin{proof}
设 \(h_j\) 如上一小节。Sobolev 范数经 Plancherel 等价于
\(\|(1+|x|)^s h_j\|_2\)；固定光滑截断的乘法在 \(H^s\) 有界，故这些范数一致受 \(CB_s\) 控制。
证明这一乘法事实，只需把乘积的 Fourier 变换写为卷积，使用
\((1+|x|)^s\le C_s(1+|x-y|)^s(1+|y|)^s\) 及 Schwartz 核的带权 \(L^1\) 范数。
同样，乘上环带中的坐标 \(\xi_\ell\) 表明
\(\|(1+|x|)^s\nabla h_j\|_2\le CB_s\)。
由 Cauchy–Schwarz，
\[
 \|h_j\|_1+\|\nabla h_j\|_1\le CB_s,\qquad
 \int_{|x|>R}|h_j(x)|\,\mathrm dx\le CB_sR^{-(s-d/2)}
 \quad(R\ge1).
\]

记 \(\gamma=s-d/2>0\)。对位移 \(v\ne0\)，低频环带 \(2^j|v|\le1\) 满足
\(\|k_j(\,\cdot-v)-k_j\|_1\le CB_s2^j|v|\)。
高频环带 \(2^j|v|>1\) 在 \(|x|>2|v|\) 上用两个尾积分，得到
\[
 \int_{|x|>2|v|}|k_j(x-v)-k_j(x)|\,\mathrm dx
 \le CB_s(2^j|v|)^{-\gamma}.
\]
两段等比级数可求和，故每个有限环带和的核及其伴随都满足一致的积分 Hörmander 条件。

\chapref{b2:ch:18}弱型证明中，坏块的估计只用该积分条件：用零平均把核减去块中心的值，再把差分积分界乘以块的 \(L^1\) 范数。好部仍用一致的 \(L^2\) 范数。因此同一证明给出有限和的一致弱型界；Marcinkiewicz 插值和伴随对偶给出一致的 \(L^p\) 界。
有限和在 \(L^2\) 中趋于 \(T_m\)，取几乎处处子列并用 Fatou，便将 \(L^p\) 界传到极限，最后用稠密性延拓。
这条证明不要求极限核有逐点梯度界，因而确实比上一版本弱。
\end{proof}
严格条件 \(s>d/2\) 用于尾积分中的正指数 \(\gamma\)。在等号处，高频环带级数不再衰减；这说明本证明失效的位置，不能把失效本身当作等号处任何可能强化的反例。

% 扩写定位：b2:subsec:220204
有限阶条件中的尺度因子不可删除。若在频率 \(\xi\) 附近的变化尺度是 \(|\xi|\)，则一次求导自然带来 \(|\xi|^{-1}\)；Mihlin 假设以 \(|\xi|^{|\alpha|}\) 恢复这个尺度。普通的“全部导数有界”在无穷频率未必提供这种衰减，也就不能直接控制重缩放后的二进核。

\ProseParagraphBreak

例如 \(m(\xi)=e^{2\pi\mathrm ia\cdot\xi}\) 的乘子是平移，故在每个 \(L^p\) 上都是等距。但其导数在无穷处不衰减，除非 \(a=0\)，不满足 Mihlin 的尺度导数条件。这个例子表明定理提供的是充分条件，而不是所有有界乘子的分类。反过来，把一个任意有界但高度不规则的符号放入 \(L^\infty\)，只能保证 \(L^2\) 有界，不能得到其他指数。

\ProseParagraphBreak

还应区分单个乘子和一族乘子的界。若 \(m_\nu\) 的 Mihlin 常数一致，则标量算子范数一致；若要在位置上先对输出求平方和，需要更强的向量论证。不过在此处统一符号条件已经足以给出这一论证：对有限族取对角乘子 \(\operatorname{diag}(m_\nu)\)，归一化后的核的算子范数由各个标量核上界的上确界控制，\(L^2(\ell^2)\) 界由逐项 Plancherel 相加得到。\cref{b2:prop:18-hilbert-valued-cz}于是给出
\[
 \left\|\left(\sum_\nu|T_{m_\nu}f_\nu|^2\right)^{1/2}\right\|_p
 \leq C_p\left\|\left(\sum_\nu|f_\nu|^2\right)^{1/2}\right\|_p.
\]
该常数与族的大小无关。若只知道各个算子的 \(L^p\) 范数一致，而不知道统一核或其他向量结构，则不能凭空推出上述平方和界。这也是前面 Bessel 平方函数识别时应核对的额外条件。
